% Materials, left half, top: the study sites (one strip: an overview and the
% five estuaries, numbered sand discs), its key and the countries in one
% caption, then what goes into MAREA and the independent checks, both as
% run-in paragraphs in body type (\normalsize); the checks carry the map's
% discs. Map slot \figSitesW x \figSitesH (391 x 54 mm, the delivered size:
% placed at 1:1); script tools/cg/fig_sites_map.py --layout strip (the strip
% draws no legend: \cgCapSites carries its key).
% \parskip is 0 down to the caption, so the figure sits under its title rule
% and the caption 3 mm under the figure, as in the other panels.
\setlength{\parskip}{0pt}%
\cgtitle{\cguse{\cgSitesTitle}{Study sites and data}}
\vspace{1.5mm}
\cgfig{sites_map_strip}{\figSitesW}{\figSitesH}{overview and the five sites
  (Villaviciosa, Ferrol, Westerschelde, Wadden Sea, Ems-Dollard) with the
  tide gauges used as checks}
\cgcap{\cguse{\cgCapSites}{Five study sites. Basemap: Natural Earth.}}
\par\vspace{4.5mm}
\setlength{\parskip}{3.2mm}%
{\term{\cguse{\cgDataTitle}{Input data}}}\
  {\color{PSlate}\cguse{\cgDataNote}{(all MAREA needs: no field surveys)}}:\
  \cguse{\cgDataBody}{\cgInputTag{1}\ Sentinel-2 images, 10\,m, 2023--2025, as
  openEO data cubes.\quad \cgInputTag{2}\ Tide at the mouth: EOT20 and
  gauges.}\par
{\term{\cguse{\cgRefTitle}{Independent checks}}}\
  {\color{PSlate}\cguse{\cgRefNote}{(never used to make the maps)}}:\
  \cguse{\cgSitesBody}{RTK-GNSS and LiDAR; tide gauges; Vaklodingen.}\par
